\section{Biochemical automata as implementations of transition systems}
The formal recognizer has two read-only cursors. A biochemical finite-state
reader instead exposes, transforms and eventually consumes a physical input
interface. To compare them professionally, specify the abstract machine
first and then state what the chemical implementation must preserve.

Let $D=(Q,\Sigma,\delta,q_0,F)$ be a deterministic finite automaton with
$\delta:Q\times\Sigma\to Q$. Extend it to words by
$\delta^*(q,\varepsilon)=q$ and
$\delta^*(q,aw)=\delta^*(\delta(q,a),w)$. It accepts $w$ exactly when
$\delta^*(q_0,w)\in F$. A missing transition in a partial specification
may be totalized with a rejecting sink. Physical suspension, however, still
needs to be recorded separately from a completed rejecting output.

As a running example, use $Q=\{E,O\}$, initial and accepting state $E$,
and transitions that preserve state on a and swap it on b. The invariant is
\begin{equation}
 \delta^*(E,w)=E\quad\Longleftrightarrow\quad |w|_b\equiv0\pmod2.
\end{equation}
It follows by induction: a adds no b, while b flips parity. For abba the
states after successive symbols are $E,O,E,E$. This is a regular-language
reader, unlike the asynchronous equal-block construction. No hidden second
cursor is provided by drawing its input as a duplex.

\subsection{A representation relation, not an analogy}
Track one designated input molecule, treating the rest of the mixture as its
physical context. This is a single-reader specification, not a claim that a
whole reaction vessel has one DFA state.
Write $I(q,aw)$ for a molecular interface intended to represent control state
$q$ and remaining input $aw$. A transition species $T_r$ represents
$r=(q,a,q')$. The physical configuration $z$ also includes free enzymes,
bound complexes, waste and other input copies. Let $\mathcal B$ be the
designated stable interfaces and define a decoding map
\[
 \pi:\mathcal B\to Q\times\Sigma^*.
\]
An intermediate enzyme-bound complex need not belong to $\mathcal B$.
This avoids assigning a new logical state after every reversible binding event.
The physical system and the logical system have different clocks.

\begin{definition}[Step refinement at stable interfaces]
Classify execution segments from $z\in\mathcal B$ to the next stable
interface $z'\in\mathcal B$ as either silent or completing. A silent segment
must satisfy $\pi(z')=\pi(z)$. A completing segment must satisfy
\[
 \pi(z)=(q,aw)\quad\Longrightarrow\quad
 \pi(z')=(\delta(q,a),w).
\]
No other stable-interface transition is allowed by this refinement contract.
A binding--dissociation excursion is silent; a commitment consumes one symbol.
Intermediate events need not have a decoded state. Terminal interfaces require
a separate output map identifying accepting and rejecting states.
\end{definition}

\Plate{07-refinement}{One logical transition can contain multiple reaction
events. The reversible association loop is silent at the logical level.
Dashed arrows decode only stable interfaces. The lower lane is a declared
reusable-transition abstraction, not the ligation-consuming mechanism of
Figure~\ref{fig:cleavage}. Completion, correct decoding and reporter formation
are separate obligations.}{fig:refinement}

\begin{proposition}[Conditional end-to-end correctness]
Suppose initial encoding is correct, every completed step refines $\delta$,
silent segments preserve the represented configuration,
each nonterminal represented input eventually completes a step, and terminal
readout reports membership in $F$ correctly. Then processing a finite word
$w$ reports exactly whether $w\in L(D)$.
\end{proposition}
\begin{proof}
After $k$ completed steps, induction gives the represented state
$\delta^*(q_0,w_1\cdots w_k)$ and unread suffix $w_{k+1}\cdots w_n$.
Progress allows the induction to reach $k=n$ rather than stall. Terminal
readout then applies the accepting-state predicate to $\delta^*(q_0,w)$.
\end{proof}
The hypothesis about progress is not implied by the step relation. A system
that binds and dissociates forever without changing its logical state can
satisfy a safety statement about completed steps while never returning an
answer. In a probabilistic model, almost-sure eventual progress is also
weaker than a deadline guarantee. Neither is equivalent to observing one
successful band in one experiment.

\subsection{Architectures have different material balances}
The 2001 ligation-based design and the later ligation-free design must not
be merged into one reaction cartoon. The 2003 report describes reversible
software--input association followed by cleavage without software--input
ligation, allowing reuse of transition molecules \cite{fuel}. This contrast
is supported here by the paper's abstract; no experimental fidelity value
or detailed sequence reconstruction is imported from it.

For theoretical analysis, a reusable-transition coarse-graining can be written
\begin{equation}
 I+T_r\ \underset{d_r}{\stackrel{k_r^+}{\rightleftharpoons}}\ C_r
 \ \xrightarrow{c_r}\ I'_r+T_r+W_r.
 \label{eq:chemical}
\end{equation}
$C_r$ is a bound intermediate and $W_r$ collects released material.
$k_r^+$ has units $\mathrm{M}^{-1}\mathrm{s}^{-1}$; $d_r,c_r$ have units
$\mathrm{s}^{-1}$. Enzyme availability is absorbed into the effective
commitment parameter $c_r$, so this is not an explicit enzyme mass balance.
It should not be used when that lumping fails.

The elementary stoichiometric changes preserve the number of transition
carriers $N_{T_r}+N_{C_r}$ in this abstraction. A consuming architecture
would omit $T_r$ from the committed products and lose that invariant. Both
require nucleotide-level balance if used as chemical models: $W_r$ cannot
simply disappear because it is absent from the logical state. Logical input
consumption is not disappearance of physical matter.

\section{Derive stochastic choice from competing pathways}
Adar and colleagues experimentally investigated competing transition species
and calibrated concentration-to-choice mappings \cite{adar}. Their results
motivate a distinction between programmed concentrations and actual transition
probabilities; they do not identify the parameters of the model below.
All rates and numerical probabilities in our reference are assigned examples.

Consider one uncommitted input with constant free transition concentrations.
Let $b_r=k_r^+[T_r]$ be its binding hazard for channel $r$. A bound complex
then either dissociates at rate $d_r$ or commits at rate $c_r$.
Assume exponential holding times, no other reactions and return to the same
free-input state after dissociation. These assumptions define a small
continuous-time Markov model; they are not automatic properties of an assay.

\subsection{One encounter and repeated attempts}
During one encounter, the probability of commitment is
\begin{equation}
 s_r=\int_0^\infty c_r e^{-(c_r+d_r)t}\,dt
     =\frac{c_r}{c_r+d_r}.
\end{equation}
Set $B=\sum_r b_r$ and suppose $B>0$ and $c_r+d_r>0$ for every active
channel. The next bound channel is $r$ with probability $b_r/B$.
The success probability of one whole attempt is
$S=\sum_r(b_r/B)s_r$. For $S>0$, unsuccessful attempts return to the same
condition, so their probabilities form a geometric series. Therefore
\begin{equation}
 p_r=\Pr(\text{eventual commitment through }r)
 =\frac{b_rs_r}{\sum_j b_js_j}.
 \label{eq:branch}
\end{equation}
This eventual branch formula is exact under the stated renewal model.
It shows why equal concentrations do not imply equal output probabilities:
association constants and commitment-versus-dissociation ratios also matter.

The expected duration $m$ to any commitment satisfies the first-step equation
\[
 m=\frac1B+\sum_r\frac{b_r}{B}\frac1{c_r+d_r}+(1-S)m.
\]
Solving it yields
\begin{equation}
 m=\frac{1+\sum_r b_r/(c_r+d_r)}{\sum_r b_rs_r}.
 \label{eq:waiting}
\end{equation}
The numerator accounts for time spent bound. Omitting it preserves eventual
branch ratios but generally predicts the wrong completion time.
The companion function \texttt{encounter\_statistics} evaluates these formulas.
Tests independently construct the transient Markov generator and solve its
linear absorption equations for both branch probabilities and mean duration.

\subsection{When an effective exponential race is justified}
If bound encounters are brief relative to waiting for new binding,
$\sum_r b_r/(c_r+d_r)\ll1$, a useful approximation treats successful
commitments as direct hazards $\lambda_r=b_rs_r$ from the free interface.
Alternatively, these hazards may be postulated as a separate empirical
coarse-grained model. Define $\Lambda=\sum_r\lambda_r$. Then
\begin{align}
 U(t)&=e^{-\Lambda t},\\
 P_r(t)&=\int_0^t\lambda_r e^{-\Lambda s}\,ds
       =\frac{\lambda_r}{\Lambda}(1-e^{-\Lambda t}).
 \label{eq:race}
\end{align}
$U$ is uncommitted probability and $P_r$ is an \emph{absolute} committed
probability, not a conditional fraction among products. Their sum is one.
If all rates vanish, $U=1$ and every $P_r=0$; branch odds are undefined,
not uniformly distributed.

\Plate{08-race}{Constant effective hazards $0.6$ and $0.4\ \mathrm{s}^{-1}$
produce two committed-output populations and an uncommitted remainder. The
curves are generated from the declared direct-race model. They must not be
read as exact finite-time predictions of the binding--dissociation model
unless its coarse-graining is justified.}{fig:race}
\Listing{race}
\begin{center}\input{\Generated/race.tex}\end{center}
At one second, correct and incorrect output probabilities are approximately
0.379272 and 0.252848, with 0.367879 uncommitted. Correctness conditional on
commitment is $0.6$, while the absolute correct-output probability is below
$0.38$. Renormalizing observed products can hide a large incomplete population.
In a pool with depleted transitions or competing inputs, rates may vary with
time and history, invalidating this simple race even if it fits a short trace.

\section{Stochastic automata and observable output}
A stochastic reader associates each symbol $a$ with a row-stochastic matrix
$P_a$ on logical states. If row vector $\alpha_0$ is the initial distribution,
then after $w=w_1\cdots w_n$,
\begin{equation}
 \alpha_n=\alpha_0P_{w_1}P_{w_2}\cdots P_{w_n},\qquad
 p_{\rm accept}=\alpha_n\boldsymbol{1}_F.
\end{equation}
Matrix multiplication sums the probabilities of alternative paths; it does
not existentially select a favorable path. A probability-valued output is not
yet a language recognizer. One can declare a cutpoint $\theta$ and accept
when $p_{\rm accept}>\theta$, but the threshold and any margin from it are
part of the specification. This chapter makes no language-class claim for
arbitrary cutpoint automata.

To retain mass honestly, extend the parity reader with absorbing lost state
$\bot$. Let $\ell$ be per-symbol loss probability and $\eta$ the probability
of flipping the intended parity among survivors. In state order $(E,O,\bot)$,
\begin{equation}
 \widetilde P_a=
 \begin{pmatrix}
 (1-\ell)(1-\eta)&(1-\ell)\eta&\ell\\
 (1-\ell)\eta&(1-\ell)(1-\eta)&\ell\\
 0&0&1
 \end{pmatrix}.
\end{equation}
For b, interchange the first two columns in the first two rows; keep the
absorbing row unchanged. This specifies a noise model, not a fit to the
2001 or 2004 experiments. A lost input is not a completed odd-parity output.
\Listing{parity_distribution}

For abba, $\eta=0.1$ and $\ell=0.05$, the generated distribution is
\begin{center}\input{\Generated/parity.tex}\end{center}
The survival probability after $n$ symbols is $(1-\ell)^n$. Conditional on
survival, let $r_n$ be the probability of correct parity. Since a noise flip
exchanges correctness and incorrectness,
\[
 r_{n+1}=(1-\eta)r_n+\eta(1-r_n).
\]
Subtract one half to get
$r_{n+1}-1/2=(1-2\eta)(r_n-1/2)$, hence
\begin{equation}
 r_n=\tfrac12[1+(1-2\eta)^n],\qquad
 p_{\rm correct}=(1-\ell)^n r_n.
\end{equation}
Two noise flips restore parity, so this is not the no-error probability
$(1-\eta)^n$. Correct output and error-free execution are distinct events.
When survival has probability zero, the conditional fraction is undefined;
the expression for absolute correct output still evaluates to zero for $n>0$.
The tests compare the recurrence with both this closed form and exhaustive
enumeration of short error/loss histories.

Finally, observation needs its own model, such as
\begin{equation}
 \mathbb E[Y_j]=b_j+\sum_{s\in\{E,O\}}H_{js}N_s.
\end{equation}
$N_s$ is output-molecule count, $H_{js}$ is expected signal in channel $j$
per molecule of state $s$, and $b_j$ is background in signal units. Identical
columns of $H$ make the two outputs unidentifiable from these measurements,
even if the logical transitions are perfect. Controls must independently
test input integrity, conversion to terminal products, state-specific
reporters and background. A professional correctness claim joins encoding,
reaction semantics, progress, stochastic error and readout evidence; it does
not replace that chain with a single successful endpoint image.
